% !PRINTED 18
\weberpage{18}
%<<CONT
turn out. But the following beautiful and important theorem
of C.\,\gesperrt{Jordan} $^{1)}$ holds:

\gesperrt{I. However the composition series of a group
$P$ may be chosen, the index series is, apart
from the arrangement, always the same.}

The proof rests on the consideration of the complementary
groups introduced in the preceding section.

Let $Q$ and $Q_1$ be normal divisors of $P$, and at the same time $Q_1$
a divisor (hence, by \S.\,4, 5. a normal divisor) of $Q$.

Then the following theorem holds:

\gesperrt{1. The group $Q/Q_1$ is a normal divisor of the
group $P/Q_1$, and the index of $Q/Q_1$ with respect
to $P/Q_1$ is equal to the index of $Q$ with respect
to $P$.}

To see its correctness, one need only consider more
closely the mode of formation of the individual groups.

Let $n, q, q_1$ be the degrees of $P, Q, Q_1$ and

\weq{3}{n = \nu q,\quad q = \nu_1 q_1,\quad n = \mu q_1,\quad \mu = \nu \nu_1.}

One now decomposes $Q$ into the subsidiary groups with respect to $Q_1$, i.e. one
sets, if $b_1, b_2 \ldots b_{\nu_1 - 1}$ are suitably determined elements in $Q$,

\weq{4}{Q = Q_1 + Q_1 b_1 + \cdots Q_1 b_{\nu_1 - 1}.}

The elements of the group $Q/Q_1$ are

\weq{5}{Q_1,\, Q_1 b_1,\, \ldots\, Q_1 b_{\nu_1 - 1}.}

If we decompose $P$ into the subsidiary groups with respect to $Q_1$, then
the elements (5) certainly occur among them, and consequently $Q/Q_1$
is a divisor of $P/Q_1$.

It therefore remains to show that this divisor is a
normal divisor.

But if we denote by $a$ any element of $P$, then
all elements of $P/Q_1$ are contained in the form $Q_1 a$, and
we therefore have only to show that, if $b$ is any element
in $Q$, every element

\weq{6}{Q_1 a^{-1} Q_1 b Q_1 a}

is contained in $Q/Q_1$, i.e. in the form $Q_1 b$.

$^{1)}$ C.\,\gesperrt{Jordan}, Traité des substitutions, Paris 1870. \gesperrt{Netto}, Substitutionentheorie,
Leipzig 1882. We follow here the proof which
\gesperrt{Hölder} has given following \gesperrt{Netto} (Math. Annalen, Bd. 34).
